%% ch08_conclusions.tex
%% Chapter 8: Conclusions
%% Thesis: Advanced Protection for Distributed Power Systems and Microgrids
%% ============================================================
\chapter{Conclusions}
\label{ch:conclusions}

\section{Summary of Contributions}
\label{sec:ch13_summary}

This thesis has developed, analysed, and validated the \textbf{Self-Adaptive
Model-Based Protection (SAMBP)} framework---a hierarchical, physics-aware
architecture for protection of future power systems dominated by
Inverter-Based Resources (IBRs). The framework treats adaptive protection
as an \emph{inverse estimation problem}: each protection element maintains
an explicit physical model of its protected device and continuously
identifies the current model parameters from real-time measurements.
Adaptation decisions are gated by a four-component confidence score and
bounded by a provably safe update law.

The thesis delivers four original contributions, each validated with
full numerical simulation results.

\subsection{C1: Adaptive Differential Protection (Chapter~\ref{ch:differential})}

The physics-constrained adaptive 87L line differential relay modulates
the bias slope $k_m^*(\kibr)$ in real time, achieving TPR\,=\,1.000 and
zero false trips across the full IBR penetration range
$\kibr \in [0.03, 0.85]$~pu. The negative-sequence element (87LN) proves
that IBR negative-sequence suppression \emph{improves} sensitivity
($\rho_2 = 2$, independent of $\kibr$). The 87T-SPRT scheme replaces
ad-hoc second-harmonic blocking with a Wald sequential probability ratio
test, providing explicit false-trip and miss-rate guarantees
($\alpha = \beta = 0.001$), and the IBR-adaptive transformer bias slope
$k_T^*(\kibr)$ closes the restrain-during-internal-fault failure mode.

\subsection{C2: Adaptive Overcurrent and Distance Protection (Chapter~\ref{ch:oc_distance})}

The SAMBP-SyncOC framework introduces inverse-estimation-based adaptive
overcurrent protection for synchronous generators. The two-pass
Levenberg--Marquardt estimator achieves $R^2 > 0.999$ with Jacobian
condition number $\kappa_n(J) < 20$ across three-phase, line-to-line,
and SLG faults. The four-component confidence-gated bounded update law
reduces spurious relay pickups by $\approx 60\%$ with no security
degradation. For distance protection, the IBR-corrected quadrilateral
characteristic achieves $P_\mathrm{detect} = 99.3\%$ and
$P_\mathrm{secure} = 99.97\%$, and the Monte Carlo reliability framework
provides the first probabilistic reliability quantification for IBR
distance protection. The microgrid protection architecture
(Chapter~\ref{ch:microgrid}) extends this to islanded operation:
ROCOF-supervised islanding detection achieves NDZ\,$< 1\%$ and 114~ms
isolation time compliant with IEEE~1547-2018, and the adaptive OC
setting-group switches within 5~ms of GOOSE receipt.

\subsection{C3: Digital Twin and Physics-Informed ML (Chapter~\ref{ch:monitoring_ml})}

The non-blocking EKF digital twin jointly estimates the state vector
$(\kibr, \Zline, \alpha)$ from IED phasor measurements, achieving
RMSE\,$\leq 0.024$~pu, $\leq 0.5\%$, $\leq 1.1\%$ respectively,
bounded by the analytically derived Cramér--Rao Lower Bound. The shadow
register pattern guarantees that EKF updates never delay primary trip
decisions. The integrated PINN fault classifier embeds three KVL/KCL
physics constraints in the loss function, achieving 96.1\% five-class
accuracy with zero physics-constraint violations. CUSUM-triggered online
learning detects fleet composition drift within 21 samples with
$< 2$~pp catastrophic forgetting.

\subsection{C4: Wide-Area Coordination and Cybersecurity (Chapter~\ref{ch:coordination_security})}

The Wide-Area Protection Coordinator integrates PMU synchrophasor data
for fault localisation (94.5\% correct-line identification at
$\sigma_V = 0.01$~pu), automatic CTI-graded relay settings dispatch
(100\% CTI-compliant in $< 3$~s), and a closed-loop
CUSUM$\to$settings$\to$GOOSE adaptive protocol completing the full
cycle in $\approx 31$~minutes with 8 of 9 steps automated.
HMAC-SHA256 authentication reduces GOOSE attack risk by 54.9\%
at 0.72~ms latency overhead, meeting IEC~61850 Class~P2 requirements.

\section{Mapping to Research Objectives}
\label{sec:ch13_objectives}

\begin{table}[htbp]
\caption{Research objectives and their fulfilment status}
\label{tab:ch13_objectives}
\small
\begin{tabularx}{\textwidth}{@{}lp{2.5cm}p{5cm}p{3cm}l@{}}
\toprule
Obj. & Chapter & Contribution & Key metric & Status \\
\midrule
O1 & \ref{ch:differential} & C1: Adaptive 87L/87T differential
  & TPR\,=\,1.000, all $\kibr$ & [P] \\
O2 & \ref{ch:oc_distance}, \ref{ch:microgrid} & C2: Adaptive OC, distance \& microgrid
  & $R^2 > 0.999$; $P_d = 99.3\%$; 114~ms & [P] \\
O3 & \ref{ch:monitoring_ml} & C3: EKF digital twin + PINN classifier
  & RMSE\,$\leq 0.024$~pu; 96.1\% & [P] \\
O4 & \ref{ch:coordination_security} & C4: WAPC + cybersecurity
  & 94.5\% localisation; 54.9\% risk reduction & [P] \\
\bottomrule
\end{tabularx}
\end{table}

\section{Limitations}
\label{sec:ch13_limitations}

\begin{enumerate}
\item \textbf{Validated scope:} All contributions are validated
      by EMT simulation and software-in-the-loop. Physical hardware-in-the-loop
      on commercial IEDs (SEL-411L) is an identified next step.

\item \textbf{Synthetic training data:} PINN training and EKF validation
      data are synthetic. Real field fault recordings from live IBR
      substations are needed for independent validation.

\item \textbf{Single-substation WAPC:} The WAPC architecture is designed
      for up to 34 buses within one substation. Multi-substation
      coordination requires a hierarchical extension.

\item \textbf{GFM model horizon:} GFM inverters are modelled as VSMs
      valid for $t \leq 150$~ms. AVR dynamics beyond this horizon are
      not captured.

\end{enumerate}

\section{Concluding Remarks}
\label{sec:ch13_concluding}

The displacement of synchronous machines by IBRs is irreversible and
accelerating. The protection crisis it creates is not a temporary
transition problem but a permanent structural feature of all future
high-IBR power systems. Legacy relay algorithms, designed for passive,
high-current, physics-determined fault current sources, fail systematically
when the fault current is controlled, software-limited, and compositionally
variable.

The SAMBP framework proposed in this thesis addresses this crisis from
first principles: by treating adaptive protection as an inverse estimation
problem, it provides a unified mathematical structure applicable to every
protection element---line differential, generator overcurrent, transformer
differential, distance, and islanding---and to every source type.
The framework is the first to integrate model identification, real-time
state estimation, and centralised coordination into a closed-loop
self-adaptive system with provable safety guarantees.

The four validated contributions establish that the framework works: relay
algorithms adapting in real-time to measured IBR fleet composition achieve
full detection coverage across the operating envelope, with quantified
reliability and security bounds. Future extensions to multi-substation
scaling, hardware-in-the-loop firmware validation, and real field fault
data will build directly on the foundational methodology, validated
algorithms, and documented performance bounds established here.
